\documentclass[12pt]{article}
\usepackage[utf8]{inputenc}
\usepackage[T1]{fontenc}
\usepackage{hyperref}
\title{Wadjabangayi Language Resources for K-12 Curriculum Development}
\author{Compiled by Claude Sonnet 5 (claude-sonnet-5) for Andr\'as Kornai}
\date{}
\begin{document}
\maketitle

\section{Introduction}
Wadjabangayi (Wadjabangai, ISO 639-3 wdy) was the language of an Aboriginal people of Queensland who lived south of Lancevale, around Maryvale and as far as Blackall. Wikipedia treats it as a dialect in the Bidjara group, and notes that the status of the dialects is not agreed. The search summary traces most information to Breen (1990) and Tindale (1974). Very little was found, and no teaching materials.

\section{Primary Educational Resources}
None found.

\section{Linguistic and Reference Materials}
\begin{itemize}
\item Wikipedia, Wadjabangai: \url{https://en.wikipedia.org/wiki/Wadjabangai}
\item Wikipedia, Bidjara language: \url{https://en.wikipedia.org/wiki/Bidjara_language}
\end{itemize}

\section{Religious and Community Resources}
A First Languages Australia news item titled ``Binbi Wadyabay'' appeared in the results. Its content was not examined, and the spelling differs from the language name, so a connection to wdy is unconfirmed: \url{https://www.firstlanguages.org.au/news/txqqske3ddlmwspc74jretljgroq1r}

\section{Digital and Audio Resources}
None found.

\section{Curriculum Development Context}
No evidence of a revival programme was found. Any work would have to start from the Bidjara-related community and from archival sources such as Breen (1990), with Traditional Owners' involvement.
\end{document}